\scanpage{198}
\mspage{7}
les données de ramification sur $S$ forment une 2-catégorie \inserted{$\unsure{\mathrm{Don\,ram}|\mathrm{Don\,rem}|\mathrm{Domram}}(S)$.} Pour $S'$ variable \inserted{dans $Et(S)$}, les 2-catégories ainsi attachées $\mathrm{Don\,ram}(S')$ forment une 2-catégorie fibrée \ldots.

Si $(\mathscr R,r)$ est défini par $Z,G$, de sorte qu'on a un objet galoisien canonique $P$ dans $\mathscr R_S$, de groupe $G$, un \deleted{foncteur} hom $(\mathscr R,r)-(\mathscr R',r')$ définit un objet $\varphi(P)$ de $\mathscr R'$, qui est galoisien \deleted{de} pour $G$, et un isomorphisme \uline{compatible avec} $G$
\[
r(\varphi(P))\overset{\alpha}{\simeq}\deleted{\unsure{r(P)}} Z
\]

\uline{Prop} 6 On trouve ainsi une équivalence entre $\underline{\mathrm{Hom}}((\mathscr R,r),(\mathscr R',r'))$ et la catégorie des couples $(P',\alpha)$, où $P'$ est un objet galoisien de $\mathscr R^{\unsure{\prime|\,}}_S$, de groupe $G$, et $\alpha$ un $G$-isomorphisme.
\[
r(P')\simeq Z
\]

[En effet, $P',\alpha$ étant donnés, on a une \uline{équivalence} $\mathscr R\cong{\mathscr R'}^{P'}$. On trouve en passant \unsure{ici|\S} :

\uline{Cor} \inserted{Pour} Un hom $(\varphi,\lambda)$ de données de ramification, \deleted{\unsure{on supp}} $\varphi$ est néc. \uline{pleinement} \uline{fidèle}. ]

En particulier

\uline{Cor} \deleted{\unsure{Cor|En}} Si $(\mathscr R',r')$ est défini par $Z',G'$, alors $\underline{\mathrm{Hom}}((\mathscr R,r),(\mathscr R',r'))$ est équivalent à la \uline{catégorie} des fibrés principaux $P'$ de $Z'$, de groupe $G$, sur lesquels $G'$ opère \deleted{\unsure{sur $P'$}} de façon compatible avec ses opérations sur $Z'$, \uline{munis} d'un isom de $G$-schémas sur $S$
\margin{$P'$ \uline{constant par morceaux} loc. (ét) sur $S$,}
% margin note written sideways
\[
P'/G'\simeq Z.
\]
% archive sheet number: 197
